% =====================================================================
\section{FreeRTOS: задачи и синхронизация}
% =====================================================================
\lessonintro
  {разделить программу на независимые задачи, организовать безопасный обмен данными, задействовать оба ядра}
  {весь проект шагов 1--6; прерывания из урока 3}
  {лампа разбита на задачи: ввод, экран, связь}
  {ничего нового из «железа»}

\subsection{Зачем нужна операционная система реального времени}

В конце урока 6 \code{loop()} совмещал пять дел, и каждое новое дело замедляло остальные.
FreeRTOS позволяет написать каждое дело как отдельную «программу» --- \textbf{задачу} --- со своим
бесконечным циклом и приоритетом. Планировщик переключает задачи каждую миллисекунду, а на
ESP32-S3 задачи могут выполняться на двух ядрах по-настоящему параллельно.

Вы уже пользовались FreeRTOS, не зная об этом: \code{setup()} и \code{loop()} выполняются внутри
задачи \code{loopTask} на ядре~1, а \code{delay()} --- это вызов FreeRTOS \code{vTaskDelay()}, который
отдаёт процессор другим задачам.

\begin{figure}[H]
\centering
\begin{tikzpicture}[state/.style={blk,circle,minimum size=19mm,fill=blkmem}]
  \node[state,fill=blkrf] (run) at (0,0) {Выпол-\\няется};
  \node[state] (ready) at (-4,0) {Готова};
  \node[state,fill=blkper] (blk) at (4,0) {Блоки-\\рована};
  \node[state,fill=blkrtc] (susp) at (0,-4.2) {Приоста-\\новлена};
  \draw[arr] (ready) to[bend left=25] node[above,font=\scriptsize]{выбрана планировщиком} (run);
  \draw[arr] (run) to[bend left=25] node[below,font=\scriptsize]{вытеснена} (ready);
  \draw[arr] (run) to[bend left=25] node[above,font=\scriptsize,align=center]{\code{vTaskDelay}, ждёт\\ очередь/мьютекс} (blk);
  \draw[arr] (blk.south) .. controls (3,-1.8) and (-3,-1.8) .. node[above,font=\scriptsize,pos=0.5,xshift=-12mm]{событие / таймаут} (ready.south);
  \draw[arr] (run) -- node[right,font=\scriptsize,pos=0.75]{\code{vTaskSuspend}} (susp);
  \draw[arr] (susp.west) to[bend left=25] node[below left,font=\scriptsize]{\code{vTaskResume}} (ready.south west);
\end{tikzpicture}
\caption{Состояния задачи FreeRTOS. Задача, ждущая в \code{vTaskDelay()}, не тратит процессорное время}
\end{figure}

\subsection{Первые задачи}

\begin{figure}[H]
\centering
\begin{tikzpicture}[x=0.9cm,y=0.9cm]
  \node[left,font=\sffamily\small] at (0,1.5) {Ядро 0};
  \node[left,font=\sffamily\small] at (0,0) {Ядро 1};
  \draw[->] (0,-0.8) -- (15,-0.8) node[right,font=\scriptsize]{$t$};
  \foreach \s/\e/\c/\l in {0/2/blkrf/Wi-Fi,2/4/blkper/b1,4/5/blkrf/Wi-Fi,5/9/gray!20/idle,9/11/blkper/b1,11/12/blkrf/Wi-Fi,12/14.5/gray!20/idle}
    \draw[fill=\c] (\s,1.2) rectangle (\e,1.8) node[midway,font=\scriptsize]{\l};
  \foreach \s/\e/\c/\l in {0/3/blkcpu/loop,3/4/blkmem/b2,4/7/blkcpu/loop,7/8/blkmem/b2,8/11/blkcpu/loop,11/12/blkmem/b2,12/14.5/blkcpu/loop}
    \draw[fill=\c] (\s,-0.3) rectangle (\e,0.3) node[midway,font=\scriptsize]{\l};
\end{tikzpicture}
\caption{Задачи на двух ядрах: пока одна ждёт, процессор занят другой}
\end{figure}

\codecap{Две независимые задачи (отдельный эксперимент, светодиоды на GPIO4 и GPIO7)}
\codeinput{l7_two_tasks}

Сравните с уроком 1: там для двух частот мигания пришлось вести отдельные таймеры на \code{millis()},
а здесь каждая задача просто «спит» столько, сколько ей нужно.

\subsection{Общие данные: очередь и мьютекс}

Когда несколько задач работают с одними данными, возникает опасность: одна задача может прочитать
структуру в тот момент, когда другая записала её лишь наполовину. Для передачи данных между
задачами используется \textbf{очередь} --- потокобезопасный FIFO-буфер.

\begin{figure}[H]
\centering
\begin{tikzpicture}
  \node[blk,fill=blkper,minimum width=3cm,minimum height=1.2cm] (p) {Задача-производитель\\ \scriptsize читает АЦП};
  \foreach \i/\v in {0/1820,1/1835,2/1790,3/,4/}{
    \node[draw,minimum width=11mm,minimum height=8mm,fill=\ifx\v\empty white\else blkmem\fi,font=\ttfamily\scriptsize]
      (c\i) at ($(p.east)+(3.0+\i*1.1,0)$) {\v};
  }
  \node[blk,fill=blkcpu,minimum width=3cm,minimum height=1.2cm,right=24mm of c4] (c) {Задача-потребитель\\ \scriptsize печатает в Serial};
  \draw[arr] (p) -- node[above,font=\ttfamily\scriptsize]{xQueueSend} (c0);
  \draw[arr] (c4) -- node[above,font=\ttfamily\scriptsize]{xQueueReceive} (c);
  \node[font=\sffamily\scriptsize,below=1mm of c2] {очередь на 5 элементов};
\end{tikzpicture}
\caption{Очередь развязывает задачи по времени}
\end{figure}

\codecap{Производитель и потребитель}
\codeinput{l7_queue}

Если же данные --- это общее \emph{состояние}, которое все читают и меняют (как \code{lamp}),
удобнее \textbf{мьютекс}: «ключ», который задача берёт перед работой с данными и возвращает после.
Пока ключ у одной задачи, остальные ждут.

\begin{table}[H]
\centering\small
\begin{tabularx}{\textwidth}{@{}l X@{}}
\toprule
Примитив & Когда использовать \\
\midrule
Очередь (\code{xQueue*}) & передача данных между задачами или из прерывания (\code{xQueueSendFromISR}) \\
Мьютекс (\code{xSemaphoreCreateMutex}) & защита общего ресурса: переменных, \code{Serial}, шины I\textsuperscript{2}C \\
Двоичный семафор & сигнал «событие произошло» (например, из прерывания в задачу) \\
Уведомления (\code{xTaskNotify}) & самый быстрый и лёгкий сигнал конкретной задаче \\
Программный таймер (\code{xTimerCreate}) & периодический вызов короткой функции без отдельной задачи \\
\bottomrule
\end{tabularx}
\caption{Основные средства синхронизации FreeRTOS}
\end{table}

\subsection{Шаг проекта 7: лампа на задачах}

\progress{7}

\begin{figure}[H]
\centering
\begin{tikzpicture}[node distance=8mm]
  \node[blk,fill=blkmem,minimum width=34mm,minimum height=14mm] (st) {\code{lamp}\\ \scriptsize + мьютекс \code{lampMutex}};
  \node[blk,fill=blkper,minimum width=30mm,minimum height=12mm,left=22mm of st] (in) {\code{inputTask}\\ \scriptsize кнопка, ручка · 10 мс\\ \scriptsize приоритет 2};
  \node[blk,fill=blkrf,minimum width=30mm,minimum height=12mm,right=22mm of st] (dsp) {\code{displayTask}\\ \scriptsize экран · 100 мс\\ \scriptsize приоритет 1};
  \node[blk,fill=blkcpu,minimum width=30mm,minimum height=12mm,below=10mm of st] (lp) {\code{loop()}\\ \scriptsize Serial, «пульс»};
  \draw[arr] (in) -- node[above,font=\scriptsize]{\code{lampToggle()}} node[below,font=\scriptsize]{\code{lampSetLevel()}} (st);
  \draw[arr] (st) -- node[above,font=\scriptsize]{\code{getLamp()}} (dsp);
  \draw[darr] (lp) -- node[right,font=\scriptsize]{\code{lampSetOn()}, \code{getLamp()}} (st);
  \node[below=1mm of dsp,font=\sffamily\scriptsize,align=center] {единственная задача,\\ работающая с I\textsuperscript{2}C};
\end{tikzpicture}
\caption{Архитектура лампы: задачи обращаются к общему состоянию только через функции под мьютексом}
\end{figure}

Здесь окупается решение из урока 4: состояние лампы меняется только через функции, поэтому защитить
его достаточно в одном месте --- внутри этих функций.

\codecap{Умная лампа, шаг 7 (изменения относительно шага 6)}
\codeinput{lamp_step7}

Кнопка теперь опрашивается каждые \SI{10}{ms} независимо от того, чем занят экран: у задачи ввода
выше приоритет, и планировщик прервёт рисование, как только ей пора работать.

\begin{caution}
Переполнение стека --- частая причина перезагрузок с сообщением \code{Stack canary watchpoint triggered}.
Проверяйте запас функцией \code{uxTaskGetStackHighWaterMark(NULL)} и увеличивайте размер стека при
использовании \code{printf}, \code{String} и больших локальных массивов.
\end{caution}

\begin{summary}
\begin{itemize}[leftmargin=*]
  \item Задача --- функция с бесконечным циклом; \code{vTaskDelay()} отдаёт процессор другим.
  \item Данные между задачами передаём очередью, общее состояние защищаем мьютексом.
  \item Каждое периферийное устройство (экран, Serial) лучше закрепить за одной задачей.
\end{itemize}
\end{summary}

\begin{tasks}
\begin{enumerate}
  \item Переведите кнопку на прерывание (урок 3): обработчик будит \code{inputTask} через \code{vTaskNotifyGiveFromISR()}.
  \item Раз в 5 секунд печатайте в Serial запас стека каждой задачи.
  \item Вынесите «пульс» в отдельную задачу и уберите \code{delay(10)} из \code{loop()}. Что изменилось?
\end{enumerate}
\end{tasks}
